\medskip

\begin{minipage}{0.5\linewidth}
Sur la figure ci-contre, les points B, A et D sont alignés. 

Les droites (BA) et (EC) sont parallèles.

\begin{enumerate}
\item Rappeler la propriété de la somme des angles d'un triangle, puis calculer la mesure de l'angle
$\widehat{\text{ACB}}$ repéré par la lettre $x$.
\item 
	\begin{enumerate}
		\item Que peut-on dire des droites (AB) et (EB) ? 

\textbf{Justifier la réponse.}
		\item En déduire la mesure de l'angle
$\widehat{\text{CBE}}$ repéré par la lettre $y$.
	\end{enumerate}
\item On s'intéresse à l'angle $\widehat{\text{ADC}}$ repéré par la lettre $z$.

\textbf{Déterminer la mesure de cet angle en expliquant chaque étape de la démarche.}
\end{enumerate}
\end{minipage}\hfill
\begin{minipage}{0.46\linewidth}
\psset{unit=0.9cm}
\begin{pspicture}(7.7,6.6)
\def\dblv{\psline(-0.05,0.1)(-0.05,-0.1)\psline(0.05,0.1)(0.05,-0.1)}
%\psgrid
\pspolygon(0.1,2.75)(1.1,0.2)(4.8,1.6)(7.2,5.8)%BECD
%\psarc(7.2,5.8){4.8}{187}{250}
\psline(2.8,3.92)(4.8,1.6)(0.1,2.75)%ACB
\uput[l](0.1,2.75){B} \uput[dl](1.1,0.2){E} \uput[dr](4.8,1.6){C} \uput[ur](7.2,5.8){D} \uput[ul]
(2.8,3.92){A} 
\rput{26}(5,4.86){\dblv}\rput{64}(6,3.7){\dblv}
\psarc(0.1,2.75){0.4}{-11}{26}\psarc(0.1,2.75){0.6}{-67}{-12}
\psarc(2.8,3.92){0.5}{201}{310}\psarc(4.8,1.6){0.5}{127}{164}
\psarc(7.2,5.8){0.6}{201}{240}
\rput(2.5,3.2){$108^\circ$}\rput(0.9,2.8){$36^\circ$}
\rput{20}(1.1,0.2){\psframe(0.25,0.25)}\rput(0.8,2.1){$y$}\rput(4.1,2){$x$}
\rput(6.5,5.2){$z$}
\end{pspicture}
\end{minipage}

\bigskip
